\documentclass[12pt,a4paper]{article}

% ---------------------------------------------------------
% Idioma y tipografía
% ---------------------------------------------------------
\usepackage[utf8]{inputenc}
\usepackage[T1]{fontenc}
\usepackage[provide=*,spanish]{babel}
\usepackage{mathpazo}
\usepackage{microtype}

% ---------------------------------------------------------
% Matemática y formato
% ---------------------------------------------------------
\usepackage{amsmath,amssymb,amsthm}
\usepackage{geometry}
\usepackage{enumitem}
\usepackage[hidelinks]{hyperref}

\geometry{
	margin=2.5cm
}

\title{Ecuaciones Diferenciales Ordinarias\\Material suplementario}
\author{Benjamín Marcolongo}
\date{}

\begin{document}

	\maketitle

	\section{Introducción}

	En las primeras clases utilizamos principalmente ecuaciones de primer orden escritas en forma normal,
	\[
	\dot{y}=f(x,y),
	\]
y buscamos soluciones expresadas como funciones explícitas \(y=\varphi(x)\). Sin embargo, ni las ecuaciones diferenciales ni sus soluciones tienen que presentarse siempre de esa manera.

	En este material utilizaremos la notación de Newton,
	\[
	\dot{y}=\frac{dy}{dx},
	\qquad
	\ddot{y}=\frac{d^2y}{dx^2}.
	\]
Aunque el punto se emplea habitualmente cuando la variable independiente es el tiempo, aquí indicará derivación respecto de $x$.

	Este material complementa la discusión anterior con tres ideas:
	\begin{enumerate}[label=\arabic*.]
		\item una solución puede quedar definida mediante una relación implícita entre \(x\) e \(y\);
		\item una EDO de primer orden puede escribirse en forma diferencial;
		\item al manipular una ecuación pueden perderse soluciones que deben estudiarse por separado.
	\end{enumerate}

	Estas cuestiones no son meramente formales. Determinan qué curvas representan realmente soluciones, en qué intervalos lo hacen y si la familia obtenida mediante un método de resolución contiene todas las soluciones posibles.

	\clearpage
	\section{Soluciones explícitas e implícitas}

	\subsection{Solución explícita}

	Una solución está dada en \textbf{forma explícita} cuando la variable dependiente aparece despejada:
	\[
	y=\varphi(x).
	\]
Por ejemplo,
	\[
	y=Ce^{3x}
	\]
es una familia explícita de soluciones de \(\dot{y}=3y\).

	\subsection{Solución implícita}

	Con frecuencia, el proceso de integración conduce a una relación de la forma
	\[
	\Phi(x,y)=C,
	\]
sin que resulte posible o conveniente despejar \(y\). En ese caso se habla de una \textbf{solución implícita}.

	Estrictamente, una solución de una EDO es una función. Por lo tanto, la relación \(\Phi(x,y)=C\) constituye una solución implícita cuando define, al menos localmente, una o más funciones diferenciables \(y=\varphi(x)\) que satisfacen la ecuación diferencial.

	Si \(\Phi\) es continuamente diferenciable y
	\[
	\frac{\partial\Phi}{\partial y}(x_0,y_0)\neq 0,
	\]
el teorema de la función implícita garantiza que, cerca de \((x_0,y_0)\), la ecuación \(\Phi(x,y)=C\) determina una única función \(y=\varphi(x)\). Derivando respecto de \(x\),
	\[
	\frac{d}{dx}\Phi\bigl(x,y(x)\bigr)=0,
	\]
y aplicando la regla de la cadena,
	\[
	\frac{\partial\Phi}{\partial x}
	+
	\frac{\partial\Phi}{\partial y}\,\dot{y}=0.
	\]
Siempre que \(\partial\Phi/\partial y\neq 0\), se obtiene
	\[
	\dot{y}
	=-\frac{\partial\Phi/\partial x}{\partial\Phi/\partial y}.
	\]
Esta expresión permite verificar directamente qué ecuación diferencial satisfacen las curvas de nivel \(\Phi(x,y)=C\).

	\subsection{Ejemplo: una familia de circunferencias}

	Consideremos
	\[
	\dot{y}=-\frac{x}{y},
	\]
definida en la región \(y\neq 0\). Separando variables,
	\[
	y\,dy=-x\,dx.
	\]
Al integrar,
	\[
	\frac{y^2}{2}=-\frac{x^2}{2}+C,
	\]
o, redefiniendo la constante,
	\[
	x^2+y^2=C.
	\]
Esta relación describe circunferencias cuando \(C>0\), pero una circunferencia completa no es la gráfica de una función \(y(x)\). Para obtener soluciones explícitas debemos considerar por separado sus dos ramas:
	\[
	y(x)=\sqrt{C-x^2}
	\qquad\text{o}\qquad
	y(x)=-\sqrt{C-x^2}.
	\]
Cada rama es solución en cualquier intervalo abierto contenido en
	\[
	(-\sqrt{C},\sqrt{C}).
	\]
Los extremos no se incluyen: allí \(y=0\), la ecuación original no está definida y la tangente a la circunferencia es vertical.

	\begin{quote}
		\textbf{Idea central.} Una curva implícita puede contener varias ramas solución y también puntos en los que deja de representar una función diferenciable de \(x\). Por eso siempre debe indicarse el intervalo de definición.
	\end{quote}

	\clearpage
	\section{Forma diferencial de una EDO}

	Una ecuación de primer orden puede escribirse en la \textbf{forma diferencial}
	\[
	M(x,y)\,dx+N(x,y)\,dy=0.
	\]
Si \(y=y(x)\), entonces \(dy=\dot{y}\,dx\). A lo largo de una curva solución,
	\[
	M(x,y)\,dx+N(x,y)\dot{y}\,dx=0,
	\]
y, por lo tanto,
	\[
	M(x,y)+N(x,y)\dot{y}=0.
	\]
En una región donde \(N(x,y)\neq 0\), esta ecuación es equivalente a la forma normal
	\[
	\dot{y}=-\frac{M(x,y)}{N(x,y)}.
	\]

	La condición \(N\neq 0\) es importante. Si \(N\) se anula, no siempre es posible despejar \(\dot{y}\), y la forma diferencial puede describir curvas con tangentes verticales que no pueden representarse localmente como \(y=y(x)\).

	\subsection{Relación con las ecuaciones separables}

	Una ecuación separable
	\[
	\dot{y}=g(x)h(y)
	\]
puede escribirse, en una región donde \(h(y)\neq 0\), como
	\[
	-g(x)\,dx+\frac{1}{h(y)}\,dy=0.
	\]
Integrar cada término produce una relación implícita:
	\[
	-\int g(x)\,dx
	+
	\int\frac{1}{h(y)}\,dy
	=C.
	\]
La notación diferencial ayuda a reconocer la separación, pero no vuelve legítima una división por cero. Los valores \(y_\ast\) que satisfacen \(h(y_\ast)=0\) deben analizarse antes de dividir por \(h(y)\).

	\subsection{Diferenciales exactos}

	Si existe una función \(\Phi(x,y)\) tal que
	\[
	M(x,y)=\frac{\partial\Phi}{\partial x},
	\qquad
	N(x,y)=\frac{\partial\Phi}{\partial y},
	\]
entonces
	\[
	M(x,y)\,dx+N(x,y)\,dy=d\Phi.
	\]
La ecuación diferencial se reduce a
	\[
	d\Phi=0,
	\]
y sus soluciones quedan definidas implícitamente por
	\[
	\Phi(x,y)=C.
	\]
	Si $M$ y $N$ poseen derivadas parciales continuas en una región abierta y simplemente conexa, la condición
	\[
	\frac{\partial M}{\partial y}
	=
	\frac{\partial N}{\partial x}
	\]
es necesaria y suficiente para que la ecuación sea exacta en esa región. La igualdad se obtiene al observar que, si $M=\partial\Phi/\partial x$ y $N=\partial\Phi/\partial y$, entonces las derivadas parciales mixtas de $\Phi$ deben coincidir.

	Por ejemplo,
	\[
	(2x+y)\,dx+(x+2y)\,dy=0
	\]
es exacta porque
	\[
	d\left(x^2+xy+y^2\right)
	=
	(2x+y)\,dx+(x+2y)\,dy.
	\]
Por lo tanto, su solución implícita es
	\[
	x^2+xy+y^2=C.
	\]
Este procedimiento será desarrollado sistemáticamente al estudiar ecuaciones exactas.

	\subsection{Equivalencia y dominio}

	Multiplicar una ecuación diferencial por una función $\mu(x,y)$ que nunca se anula no modifica sus curvas solución. En cambio, si $\mu$ puede valer cero, la nueva ecuación puede admitir curvas adicionales contenidas en el conjunto $\mu(x,y)=0$. Recíprocamente, dividir por $\mu$ elimina cualquier solución que se encuentre en ese conjunto.

	Por lo tanto, las expresiones
	\[
	M\,dx+N\,dy=0
	\]
y
	\[
	\mu M\,dx+\mu N\,dy=0
	\]
son equivalentes solamente en una región donde $\mu\neq 0$. Toda transformación algebraica de una EDO debe ir acompañada por una revisión de su dominio y de los factores que fueron cancelados.

	\clearpage
	\section{Soluciones singulares}

	Al resolver una EDO suele obtenerse una familia uniparamétrica de soluciones,
	\[
	\Phi(x,y,C)=0.
	\]
Una \textbf{solución singular} es una solución que no puede obtenerse asignando ningún valor a \(C\) en la familia encontrada.

	No debe confundirse una solución singular con una solución particular. Una solución particular se obtiene fijando el parámetro \(C\); una solución singular queda fuera de toda la familia.

	\subsection{Ejemplo: una solución perdida al dividir}

	Consideremos la ecuación
	\[
	\dot{y}=x\sqrt{y},
	\qquad y\geq 0.
	\]
La función
	\[
	y(x)=0
	\]
es una solución, ya que ambos miembros de la ecuación son nulos.

	Para buscar soluciones positivas dividimos por \(\sqrt{y}\):
	\[
	\frac{dy}{\sqrt{y}}=x\,dx.
	\]
Integrando,
	\[
	2\sqrt{y}=\frac{x^2}{2}+C_1,
	\]
y redefiniendo la constante,
	\[
	\sqrt{y}=\frac{x^2}{4}+C.
	\]
En los intervalos donde el miembro derecho es no negativo,
	\[
	y(x)=\left(\frac{x^2}{4}+C\right)^2.
	\]
En particular, si se buscan soluciones definidas para todo \(x\in\mathbb{R}\), debe tomarse \(C\geq 0\). Para \(C=0\) se obtiene
	\[
	y(x)=\frac{x^4}{16},
	\]
pero ningún valor constante de \(C\) produce la solución idénticamente nula. En el sentido de la definición anterior, \(y=0\) es una solución singular respecto de esta familia.

	La solución se perdió exactamente en el paso en el que dividimos por \(\sqrt{y}\). Este es el mismo riesgo señalado en la Clase II al resolver ecuaciones separables.

	\subsection{Soluciones singulares como envolventes}

	En la teoría clásica, el ejemplo característico de solución singular aparece como la envolvente de una familia de soluciones. Consideremos la ecuación de Clairaut
	\[
	y=x\dot{y}+\dot{y}^{2}.
	\]
La familia de rectas
	\[
	y=Cx+C^2
	\]
satisface la ecuación para todo \(C\in\mathbb{R}\). Derivando la ecuación original,
	\[
	\dot{y}
	=
	\dot{y}
	+x\ddot{y}
	+2\dot{y}\ddot{y},
	\]
de donde
	\[
	\left(x+2\dot{y}\right)\ddot{y}=0.
	\]
La rama \(\ddot{y}=0\) produce la familia de rectas. La otra posibilidad es
	\[
	\dot{y}=-\frac{x}{2}.
	\]
Al sustituirla en la ecuación original se obtiene
	\[
	y_s(x)=-\frac{x^2}{4}.
	\]
Esta parábola también satisface la EDO, pero no coincide con ninguna recta de la familia. Además, es tangente a cada una de ellas para un valor apropiado de \(x\); por eso constituye su envolvente.

	\subsection{Una precisión terminológica}

	El término «solución singular» no se usa de manera completamente uniforme en todos los textos. En un sentido amplio, se aplica a una solución que no pertenece a la familia general obtenida. En un sentido más restringido, se reserva para soluciones envolventes, como la parábola del ejemplo de Clairaut.

	En este curso utilizaremos siempre una formulación explícita: indicaremos si una solución fue excluida por una división o si es una verdadera envolvente de una familia. Esto evita que la terminología oculte el mecanismo matemático.

	\begin{quote}
		\textbf{Regla práctica.} Antes de dividir por una expresión que depende de \(x\), de \(y\) o de sus derivadas, iguale esa expresión a cero y compruebe por separado si aparecen soluciones adicionales.
	\end{quote}

	\clearpage
	\section{Resumen}

	\begin{itemize}[leftmargin=*]
		\item Una relación \(\Phi(x,y)=C\) representa una solución implícita cuando define localmente una función diferenciable que satisface la EDO.
		\item Una curva implícita completa no tiene por qué ser la gráfica de una única función \(y(x)\); puede ser necesario separarla en ramas.
		\item La forma diferencial \(M\,dx+N\,dy=0\) equivale a \(\dot{y}=-M/N\) solamente donde \(N\neq 0\).
		\item Si \(M\,dx+N\,dy=d\Phi\), las soluciones están dadas por las curvas de nivel \(\Phi(x,y)=C\).
		\item Dividir por una expresión que puede anularse puede eliminar soluciones.
		\item Una solución singular satisface la EDO, pero no pertenece a la familia uniparamétrica considerada. En los ejemplos clásicos puede aparecer como la envolvente de esa familia.
	\end{itemize}

	\section*{Bibliografía}

	\begin{itemize}[leftmargin=*]
		\item Zill, Dennis G. \textit{A First Course in Differential Equations with Modeling Applications}. 12.\textsuperscript{a} edición, edición métrica internacional. Blue Kingfisher, 2023. Capítulos 1 y 2. ISBN: 9798214038209.
	\end{itemize}

	\subsection*{Nota sobre la elaboración del material}

	Este material fue elaborado por Benjamín Marcolongo con la asistencia de \textbf{GPT-5.6 Sol}, un modelo de OpenAI utilizado a través del entorno Codex, para la organización, redacción y revisión del contenido.

\end{document}
